%%%%%%%%%%%%%%%%%%%%%%%%%%%%%%%%%%%%%%%%%%%%%%%%%%%%%%%%%%%%%%%%%%%%%%%%%%%
\chapter{Ranking, BM25, and RRF}
\label{ch:ranking}
%%%%%%%%%%%%%%%%%%%%%%%%%%%%%%%%%%%%%%%%%%%%%%%%%%%%%%%%%%%%%%%%%%%%%%%%%%%

%%%%%%%%%%%%%%%%%%%%%%%%%%%%%%%%%%%%%%%%%%%%%%%%%%%%%%%%%%%%%%%%%%%%%%%%%%%
\section{BM25 --- Best Match 25}
%%%%%%%%%%%%%%%%%%%%%%%%%%%%%%%%%%%%%%%%%%%%%%%%%%%%%%%%%%%%%%%%%%%%%%%%%%%

BM25 is the dominant full-text ranking algorithm used in search engines (Elasticsearch,
Lucene, Solr, Tantivy). It scores documents against a query based on term frequency and
inverse document frequency with saturation corrections:

\[
\text{score}(D, Q) = \sum_{t \in Q} \text{IDF}(t) \cdot \frac{\text{tf}(t,D) \cdot (k_1 + 1)}
  {\text{tf}(t,D) + k_1 \cdot \left(1 - b + b \cdot \frac{|D|}{\text{avgdl}}\right)}
\]

Where:

\begin{itemize}
    \item $\text{tf}(t, D)$ --- how many times term $t$ appears in document $D$
    \item $\text{IDF}(t)$ --- how rare the term is across all documents (log-scaled)
    \item $k_1$ --- term frequency saturation (typically 1.2--2.0); prevents a term appearing
          100$\times$ from dominating
    \item $b$ --- length normalization (typically 0.75); penalizes long documents
    \item $\text{avgdl}$ --- average document length in the corpus
\end{itemize}

\textbf{Why BM25 matters for RAG:} Vector search finds semantically similar content but can
miss exact keyword matches: names, version numbers, code identifiers, rare technical terms.
BM25 excels at exact match and rare term recall. Hybrid search combines both.

\textbf{Implementation notes:} memvid uses Tantivy's BM25. MongoDB Atlas Search uses Lucene's.
Qdrant uses FastEmbed's sparse encoding which approximates BM25.

%%%%%%%%%%%%%%%%%%%%%%%%%%%%%%%%%%%%%%%%%%%%%%%%%%%%%%%%%%%%%%%%%%%%%%%%%%%
\section{RRF --- Reciprocal Rank Fusion}
%%%%%%%%%%%%%%%%%%%%%%%%%%%%%%%%%%%%%%%%%%%%%%%%%%%%%%%%%%%%%%%%%%%%%%%%%%%

When you run both BM25 and vector search on the same query, you get two ranked lists. The
lists contain overlapping but not identical results with incomparable scores (BM25 scores are
not in the same range as cosine similarity values).

RRF merges them by rank position, not by raw score:

\[
\text{RRF\_score}(d) = \sum_{L} \frac{1}{k + \text{rank}_L(d)}
\]

Where $k$ is a constant (typically 60). The $k$ value dampens the advantage of rank~1 over
rank~2---it prevents a single list from dominating just because one result was ranked first.

\textbf{Example with $k=60$ and two lists:}

\begin{Verbatim}
doc A: rank 1 in BM25, rank 4 in vector  → 1/61 + 1/64 = 0.0321
doc B: rank 2 in BM25, rank 2 in vector  → 1/62 + 1/62 = 0.0323
doc C: rank 1 in vector, not in BM25     → 0    + 1/61 = 0.0164
\end{Verbatim}

Doc~B wins because consistent top performance across both lists beats a single dominant rank.
Doc~C is penalized for appearing in only one list.

RRF requires no calibration of score scales across different retrieval systems---only ranks
matter. This makes it robust for combining any retrieval methods without normalization.

memvid's \texttt{ask()} method uses RRF with $k=60$ as the default fusion strategy.

%%%%%%%%%%%%%%%%%%%%%%%%%%%%%%%%%%%%%%%%%%%%%%%%%%%%%%%%%%%%%%%%%%%%%%%%%%%
\section{nDCG in Practice: Using Rankings to Make Better Decisions}
%%%%%%%%%%%%%%%%%%%%%%%%%%%%%%%%%%%%%%%%%%%%%%%%%%%%%%%%%%%%%%%%%%%%%%%%%%%

Understanding nDCG at a practical level helps you make better decisions about your retrieval
stack. The key insight is that nDCG measures ranking quality, not just membership.

\subsection{Why nDCG Is Better Than Plain Recall for Ranking Quality}

Recall asks: did the system retrieve relevant items? It does not strongly care whether the
best item was rank~1 or rank~10. nDCG does care. This is why:

\begin{itemize}
    \item Recall is often more important for stage-1 candidate generation
    \item nDCG is often more important for top-of-list quality and reranking
\end{itemize}

\subsection{Why \texttt{@10} Is Common}

\texttt{nDCG@10} is common because many applications only care about the first screenful of
results or the first prompt window. A RAG pipeline may pass only the top 5 to top 10 chunks
into generation. If your system only ever uses the top 5 chunks, then \texttt{nDCG@5} may be
even more informative for your internal evals.

\subsection{nDCG with Binary vs Graded Judgments}

With binary judgments, nDCG mostly reflects whether relevant items are ranked early.
With graded judgments, nDCG becomes more expressive---it can reward systems that rank the
very best documents above merely acceptable ones. This makes graded nDCG especially valuable
for reranking benchmarks.

\subsection{What nDCG Does Not Tell You}

nDCG does not directly tell you:

\begin{itemize}
    \item Whether your retriever found all relevant items in the corpus
    \item How expensive the retrieval system is
    \item Whether the retrieved text is the best chunking strategy for generation
    \item Whether your model generalizes outside the benchmark
\end{itemize}

That is why nDCG should usually be read together with Recall@k and domain-specific benchmark
slices or internal evaluation.

%%%%%%%%%%%%%%%%%%%%%%%%%%%%%%%%%%%%%%%%%%%%%%%%%%%%%%%%%%%%%%%%%%%%%%%%%%%
\section{SPO and Entity Graphs}
%%%%%%%%%%%%%%%%%%%%%%%%%%%%%%%%%%%%%%%%%%%%%%%%%%%%%%%%%%%%%%%%%%%%%%%%%%%

\textbf{SPO triplets} are a way to represent facts as structured relationships:

\begin{Verbatim}
Subject        Predicate       Object
───────────    ─────────────   ──────────────
Alice          works_at        Acme Corp
Bob            lives_in        San Francisco
project-X      depends_on      library-Y
\end{Verbatim}

SPO extraction parses free text and produces a graph of entities and their relationships.
In retrieval systems, this enables queries that pure vector search cannot handle well:

\begin{itemize}
    \item ``who works at Acme Corp'' --- entity lookup, not semantic similarity
    \item ``what does project-X depend on'' --- graph traversal
    \item ``find all relationships involving Alice'' --- entity-centric search
\end{itemize}

\textbf{MemoryCards} (as implemented in memvid) are typed records built on top of SPO:

\begin{description}
    \item[Fact] A general assertion (``Paris is the capital of France'')
    \item[Preference] A user preference (``user prefers dark mode'')
    \item[Event] A timestamped occurrence (``meeting with Alice on 2026-03-15'')
    \item[Relationship] A link between entities (``Alice reports to Bob'')
\end{description}

In agent memory systems, the SPO graph complements vector search: vector search finds
semantically similar content; graph search finds structured facts about known entities.

%%%%%%%%%%%%%%%%%%%%%%%%%%%%%%%%%%%%%%%%%%%%%%%%%%%%%%%%%%%%%%%%%%%%%%%%%%%
\section{Choosing the Right Retrieval Mix}
%%%%%%%%%%%%%%%%%%%%%%%%%%%%%%%%%%%%%%%%%%%%%%%%%%%%%%%%%%%%%%%%%%%%%%%%%%%

\begin{center}
\begin{tabular}{ll}
\toprule
\textbf{Query type} & \textbf{Best retrieval method} \\
\midrule
Semantic / conceptual & Dense vector search \\
Exact keyword (names, IDs, version numbers) & BM25 \\
Mixed / general & Hybrid (BM25 + vector with RRF) \\
Multi-aspect complex queries & ColBERT or reranker \\
Structured entity queries & SPO graph search \\
\bottomrule
\end{tabular}
\end{center}

In practice, running hybrid retrieval (BM25 + vector) as the default and adding a reranker
for final ordering captures the vast majority of production retrieval improvements. SPO graph
search is valuable for agent memory and knowledge graph use cases but requires additional
infrastructure.
